% Fixed_Points_I.tex
% Requires: series.sty in your TEXINPUTS (same folder is fine).
\documentclass[11pt]{article}
\usepackage{series}

% --- Paper metadata (override any defaults from series.sty, if needed) ---
\hypersetup{
  unicode=true,
  pdftitle={Fixed Points I: Fixed Points over Multi--Bundle Manifolds},
  pdfauthor={Peter Nero}
}

% --- Local macros (paper-specific, minimal) ---
% NOTE: use \providecommand to avoid collisions with series.sty (shared macros live there).
\providecommand{\Bint}{B_{\mathrm{int}}}
\providecommand{\PiCoh}{\Pi_{\mathrm{coh}}}
\providecommand{\Hcoh}{H_{\mathrm{coh}}}
\providecommand{\Qop}{Q}
\providecommand{\Aop}{A}
\providecommand{\Aeps}{A_\varepsilon}
\providecommand{\Lip}{L_{\mathrm{Lip}}} % Lipschitz constant (series-wide)
\providecommand{\Ran}{\mathrm{Ran}}
\providecommand{\dom}{\mathrm{dom}}
\providecommand{\spec}{\sigma}

% PDF-safe strings in bookmarks (fallbacks if series.sty doesn't define them)
\providecommand{\FPpdf}[2]{\texorpdfstring{#1}{#2}}
\providecommand{\lamstar}{\ensuremath{\lambda^{\ast}}}
\providecommand{\lamone}{\ensuremath{\lambda_{1}}}
\providecommand{\epsdown}{\ensuremath{\varepsilon\downarrow 0}}
\providecommand{\Sone}{\ensuremath{S^{1}}}



\title{Fixed Points I: Fixed Points over Multi--Bundle Manifolds\\
\large A conditional functional-analytic framework with spectral-bundle projections}
\author{Peter Nero}
\date{July 2026}

\begin{document}
\maketitle
\seriestagline

\begin{abstract}
We develop a conditional fixed-point framework for gradient flows on internal fiber
products of bundles over a base manifold $Y$ of dimension $d := \dim Y$. Given smooth
bundles $\pi_n : B_n \to Y$ with compact fibers $F_n$ of uniformly bounded geometry, we
construct spectral-bundle projections onto the joint harmonic sector and study gradient
flows generated by coercive analytic energies $C : H^1(\Bint)\to \bbR$ on the internal product
\[
\Bint := B_1 \times_Y \cdots \times_Y B_k .
\]
We establish from first principles:
(i) precise geometry and anisotropic Sobolev structure;
(ii) uniform $H^1$ bounds for the coherent projector via Riesz--spectral calculus under bounded geometry;
(iii) analytic-semigroup well-posedness and smoothing for $\Phi_t$;
(iv) existence of projected fixed points by a variational method on the coherent sector; alternatively by Schauder (compact base with base-regularized smoothing) or by Darbo using the Kuratowski measure under a confining base potential;
(v) uniqueness either from strong convexity on the coherent sector or via an inertial-manifold reduction under a clean spectral gap condition;
(vi) quantitative Galerkin approximation for the uniformly elliptic steady problem with C\'ea-type quasi-optimal error;
(vii) examples (with explicit constants) and \L{}ojasiewicz--Simon convergence.
All results are analytic/control statements under the displayed hypotheses. The
Riemannian base is not automatically physical Lorentzian spacetime, the flow
parameter is not automatically physical time, and the abstract fiber product
does not by itself select the canonical MTT internal topology.
\end{abstract}

\section*{Revision note for this edition}
\addcontentsline{toc}{section}{Revision note for this edition}
\begin{description}
\item[Supersedes.] \emph{Fixed Points I: Fixed Points over Multi--Bundle
Manifolds}, version~6.
\item[Reason.] The version~6 technical correction was published before its
paper-specific reader map, theorem-cluster explanations, and final boundary
discussion were incorporated into the released PDF.
\item[Resolution.] Version~7 publishes the reviewed explanatory source.  It
adds a plain-language model, a six-stage argument map, interpretation around
the principal theorem clusters, and a conclusion that separates analytic
control from physical realization.  No theorem is strengthened by this
editorial repair.
\item[Retained result.] Conditional projected fixed-point existence,
coherent-sector uniqueness routes, and convergence remain valid under the
displayed analytic hypotheses.
\item[Remaining boundary.] The abstract Riemannian control geometry, flow
parameter, and fiber product do not select physical spacetime or the q79
internal carrier.
\end{description}

\section*{Prior technical revision (version 6)}
\addcontentsline{toc}{section}{Prior technical revision (version 6)}
Version~6 repaired the spectral-calculus and semigroup estimates, stated
valid Schauder and Darbo routes, corrected the confining-space norm, and
proved a Cea-type Galerkin estimate for the elliptic steady problem.  Those
technical changes are retained unchanged in this edition; their detailed
audit remains part of the source archive.

\section*{How to read this paper}
\addcontentsline{toc}{section}{How to read this paper}
This paper answers a deliberately limited question: once a concrete model has
supplied a family of compact internal fibers, a coherent projector, and an
energy, what analytic conditions are sufficient to guarantee that the
projected dynamics has a genuine equilibrium?  It does not derive those
geometric inputs from MTT.  It develops the functional-analytic machine that
can be applied after they have been selected.

\paragraph{The central picture in plain language.}
A field on the total bundle contains two kinds of variation.  Its coherent
part lies in the fiberwise harmonic modes retained by $\PiCoh$; its
incoherent part lies in the positive-spectrum complement.  The fiber
Laplacian damps the latter, while the nonlinear term determines how the
retained modes interact.  A projected fixed point is physically useful only
if projection has not hidden continuing motion.  Coherence invariance and
energy dissipation are the conditions that turn such a projected return into
a true stationary state.

\paragraph{Argument map.}
The proof architecture has six stages.
\begin{enumerate}[label=\textup{\arabic*.}]
\item Sections~1--2 define the fiber product and use the uniform fiber gap to
construct a bounded harmonic projector.
\item Section~3 places a dissipative semilinear flow on that space and proves
the key implication from projected fixed point to actual equilibrium.
\item Section~4 gives three alternative existence routes.  The variational
route is the shortest; Schauder is useful on compact bases with smoothing;
Darbo replaces global compactness by confinement and tail control on
noncompact bases.
\item Section~5 separates existence from uniqueness.  Strong monotonicity is
the direct route, while an inertial manifold provides a finite-dimensional
route when global convexity is unavailable.
\item Section~6 shows how a Galerkin computation approximates the selected
steady state once a coercive elliptic problem has been specified.
\item Sections~7--10 test the hypotheses in examples, establish asymptotic
convergence, and explain the controlled removal of auxiliary base
regularization.
\end{enumerate}
These routes are alternatives, not a list of conclusions available under one
minimal assumption set.  In particular, an existence theorem does not supply
uniqueness or a rate, and a Galerkin estimate does not select the geometric
operator being discretized.

\paragraph{What should be carried forward.}
Later Fixed Points papers may use the projector, equilibrium, and
approximation results by citing the corresponding statements here and
checking their hypotheses.  They should not restate these theorems, nor infer
Lorentzian time, a physical compactification, or a Standard Model vacuum from
them.  Those are separate realization and identification problems.

\tableofcontents

% ------------------------------------------------------------
\section*{Standing Series Axiom (shared across the FP series)}
\addcontentsline{toc}{section}{Standing Series Axiom (shared across the FP series)}
\begin{assumption}[Series Axiom: bounded geometry + uniform spectral gap + coherent projection]
\leavevmode
\begin{enumerate}[label=\textup{(SA\arabic*)}]
\item \textbf{Base geometry.} $(Y,g_Y)$ is a smooth, complete Riemannian manifold. When uniform constants involving base derivatives are required, we additionally assume $Y$ has bounded geometry (uniform bounds on curvature and injectivity radius on the region of interest).
\item \textbf{Fiber bounded geometry.} Each bundle $\pi_n:B_n\to Y$ has compact fibers $F_n(y)$ whose Riemannian geometry is uniformly bounded in $y$ (Assumption~\ref{ass:bounded-geom} below).
\item \textbf{Uniform fiber spectral gap.} The nonnegative fiber Laplacians $\Delta_{F_n(y)}$ have a uniform first positive eigenvalue bound $\lamstar>0$ (Assumption~\ref{ass:gap}).
\item \textbf{Coherent projector.} The fiberwise Riesz projector $\PiCoh$ onto the joint harmonic sector is well-defined and bounded on the Sobolev scales used (Section~\ref{sec:projectors}).
\item \textbf{Coherence invariance (when invoked).} For equilibrium identification and coherent-sector variational arguments we assume $N(\Ran \PiCoh)\subset \Ran\PiCoh$ (Assumption~\ref{ass:coh-inv}).
\end{enumerate}
\end{assumption}

% ------------------------------------------------------------
\section{Geometry of the internal product and standing hypotheses}

The first task is to separate what belongs to geometry from what belongs to
dynamics.  The base labels the family of internal fibers; the fiber product
places all selected internal structures over the same base point.  None of
this yet says how a configuration evolves.  It only provides the space on
which the projector and the energy will act.

\subsection{Internal fiber product}
Let $Y$ be a smooth $d$-dimensional manifold with complete Riemannian metric $g_Y$.
Let $\pi_n : B_n \to Y$ be smooth bundles with compact fibers $F_n(y)$.

\begin{remark}[Control geometry versus physical spacetime]
The base $(Y,g_Y)$ is the Riemannian geometry used by the elliptic and parabolic
estimates in this paper. Identifying it with a Cauchy slice, Euclideanized
background, or control geometry associated with a Lorentzian spacetime requires
a separate realization theorem. No Lorentzian causal or dynamical conclusion
is inferred here.
\end{remark}

The \emph{internal bundle product} (fiber product over $Y$) is
\[
\Bint := B_1 \times_Y \cdots \times_Y B_k
:= \{(b_1,\dots,b_k): \pi_1(b_1)=\cdots=\pi_k(b_k)\in Y\},
\]
with projection $\pi:\Bint\to Y$, $\pi(b_1,\dots,b_k)=y$.

\paragraph{Universal property.}
$\Bint$ comes with projections $\mathrm{pr}_n:\Bint\to B_n$ such that $\pi_n\circ \mathrm{pr}_n=\pi_m\circ \mathrm{pr}_m$ for all $n,m$.
For any manifold $M$ with maps $f_n:M\to B_n$ satisfying $\pi_1\circ f_1=\cdots=\pi_k\circ f_k$,
there is a unique $f:M\to\Bint$ with $\mathrm{pr}_n\circ f=f_n$.

We equip each fiber
\[
F(y):=\pi^{-1}(y)\cong F_1(y)\times \cdots \times F_k(y)
\]
with a smooth product-type Riemannian metric $g_{F(y)}$ varying smoothly with $y$
under bounded-geometry assumptions below, and equip $\Bint$ with the induced product-type metric.

\begin{remark}[Scope of the fiber-product model]
The displayed product is an abstract analytic construction. In an MTT
realization with overlapping vertical structures, or nested structures after
explicit inclusion maps have been supplied---in particular when there is a
shared central circle---one must work on the selected six-dimensional internal
space with the corresponding joint vertical operators; the shared structure
must not be counted as an additional independent factor. The theorems below
apply after that model supplies the required bounded geometry, joint projector,
and gap hypotheses.
\end{remark}

\subsection{Bounded geometry and uniform spectral gap}

\begin{assumption}[Bounded geometry of fibers]\label{ass:bounded-geom}
There exist $K_0,\iota_0>0$ and an integer $m\ge 2$ such that for all $y\in Y$ and each $n=1,\dots,k$,
the fiber $(F_n(y),g_{F_n(y)})$ has $|\sec|\le K_0$, $\mathrm{injrad}\ge \iota_0$,
and the fiber metric depends $C^m$-smoothly on $y$ with all $y$-derivatives up to order $m$
bounded uniformly on $Y$ (in bounded-geometry coordinate charts on the base, when needed).
\end{assumption}

\paragraph{Nonnegative Laplacians.}
Let $\Delta_{F_n(y)}$ denote the \emph{nonnegative} Laplace--Beltrami operator on $F_n(y)$ (so $\Delta\ge 0$ and $\ker\Delta$ are the fiberwise harmonic modes).
Define the fiber operator on the product fiber $F(y)$ by
\[
\Delta_F(y) := \sum_{n=1}^k \Delta_{F_n(y)}.
\]

\begin{assumption}[Uniform fiber spectral gap]\label{ass:gap}
There exists $\lamstar>0$ such that
\[
\lamstar \le \inf_{y\in Y}\ \min_{1\le n\le k}\ \lambda_1\!\big(\Delta_{F_n(y)}\big),
\]
where $\lambda_1(\Delta_{F_n(y)})$ is the first \emph{positive} eigenvalue of the nonnegative fiber Laplacian.
Moreover, $\dim\ker\Delta_{F_n(y)}$ is constant in $y$ for each $n$ (hence the harmonic bundle ranks are constant).
\end{assumption}

\subsection{Function spaces and anisotropic norms}

Let $L^2 := L^2(\Bint)$ with the product volume measure. Let $\nabla_F$ denote the gradient along fibers.
Define the anisotropic (fiber) Sobolev space
\[
H^1_F := \{\Psi\in L^2:\ \nabla_F\Psi\in L^2\},
\qquad
\|\Psi\|_{H^1_F}^2 := \|\Psi\|_{L^2}^2 + \|\nabla_F\Psi\|_{L^2}^2.
\]
We use $H_0 := L^2$ and $H_F := H^1_F$ as the primary Banach spaces for mapping properties.

\begin{definition}[Full Sobolev space]
We write
\[
H^1(\Bint)
:=
\big\{\Psi\in L^2(\Bint):\ \nabla_Y\Psi\in L^2,\ \nabla_F\Psi\in L^2\big\},
\]
with norm
\[
\|\Psi\|_{H^1}^2 := \|\Psi\|_{L^2}^2 + \|\nabla_Y\Psi\|_{L^2}^2 + \|\nabla_F\Psi\|_{L^2}^2.
\]
\end{definition}

\begin{remark}[Why $H^1_F$ is the correct space for fiber coercivity]
The fiber operator $\Delta_F$ controls $\|\nabla_F\Psi\|_{L^2}$ but does not control $\|\nabla_Y\Psi\|_{L^2}$.
Thus a Poincar\'e-type inequality driven solely by $\Delta_F$ fails on the full $H^1(\Bint)$, while it holds naturally in $H^1_F$.
This is the key coercivity correction used throughout the series.
\end{remark}

\paragraph{Coherent sector notation.}
Let $\PiCoh$ be the coherent projector (defined in Section~\ref{sec:projectors}) and set
\[
\Hcoh := \Ran(\PiCoh)\cap H^1(\Bint),\qquad \Qop := \Id - \PiCoh.
\]

% ------------------------------------------------------------
\section{Projectors via fiberwise spectral calculus}\label{sec:projectors}

The spectral gap now does the decisive organizational work.  It separates
the zero modes continuously from every positive fiber mode, so the coherent
sector is not chosen by an informal truncation: it is the range of a Riesz
projector.  Uniformity of the gap is what permits one projector bound to hold
over the whole base rather than only fiber by fiber.

For each $y\in Y$ define the self-adjoint nonnegative operator
\[
\Aop(y) := \Delta_F(y)\quad \text{on } L^2(F(y)),\qquad \dom(\Aop(y))=H^2(F(y)).
\]
By Assumption~\ref{ass:gap}, $\spec(\Aop(y))\cap (0,\lamstar)\!=\!\varnothing$ uniformly in $y$.

Let $\Gamma$ be the positively oriented circle of radius $\lamstar/2$ centered at $0$.
Define the Riesz projector onto $\ker \Aop(y)$ by
\[
\Pi(y) := \frac{1}{2\pi i}\int_{\Gamma} (z-\Aop(y))^{-1}\,dz.
\]
Define $\PiCoh$ on $L^2(\Bint)$ fiberwise by
\[
(\PiCoh\Psi)(\cdot,y) := \Pi(y)\Psi(\cdot,y).
\]

\begin{lemma}[Uniform boundedness of $\PiCoh$ on $L^2$ and $H^1_F$]
Under Assumptions~\ref{ass:bounded-geom}--\ref{ass:gap}, there exists $C_\Pi\ge 1$ such that
\[
\|\PiCoh\|_{L^2\to L^2}\le 1,\qquad \|\PiCoh\|_{H^1_F\to H^1_F}\le C_\Pi.
\]
In the $y$-independent flat case one may take $C_\Pi=1$.
\end{lemma}

\begin{remark}[On $\Ran(\PiCoh)$ for scalar fields]
If the fibers are connected and $\PiCoh$ projects scalar functions onto $\ker\Delta_F$, then $\Ran(\PiCoh)$ consists of fiberwise constants and $\nabla_F(\PiCoh\Psi)=0$, so $\|\PiCoh\Psi\|_{H^1_F}=\|\PiCoh\Psi\|_{L^2}\le \|\Psi\|_{L^2}$ (hence $C_\Pi=1$).
The general $C_\Pi\ge 1$ form is kept for vector-valued settings (e.g.\ bundles of forms/spinors) where $\ker\Delta$ may have rank $>1$.
\end{remark}

\begin{lemma}[Boundedness on full $H^1$]\label{lem:PiH1}
Under Assumptions~\ref{ass:bounded-geom}--\ref{ass:gap} (and base bounded-geometry control when $y$-derivatives are invoked),
$\PiCoh$ extends to a bounded map $H^1\to H^1$, hence $\Ran(\PiCoh)$ is closed in $H^1$.
\end{lemma}

\begin{remark}[Proof references for spectral projectors]
Uniform control of the Riesz--spectral projector on Sobolev scales follows from the Riesz--Dunford formula together with parameter-dependent elliptic regularity on bounded-geometry families; see, for example, \cite{Kato1976,Hebey2000}.
\end{remark}

\begin{remark}[Automatic commutation on the domain]
Since $\Pi(y)$ projects onto $\ker \Aop(y)$, we have $\Aop(y)\Pi(y)=0$ and $\Pi(y)\Aop(y)=0$ on $\dom(\Aop(y))$.
Thus $[\Aop,\PiCoh]=0$ on $\dom(\Aop)$, and by density/continuity on $L^2$.
Since $\PiCoh$ is a Riesz projector defined via the functional calculus of
$\Aop$, it commutes with $\Aop$ and hence with $e^{-t\Aop}$.
\end{remark}

% ------------------------------------------------------------
\section{Gradient flow: well-posedness and smoothing}

\subsection{Energy, decomposition, and basic assumptions}

Let $C:H^1(\Bint)\to \bbR$ be a real-analytic functional such that its $L^2$-gradient decomposes as
\[
\nabla C(\Psi) = \Aop\Psi + N(\Psi),
\]
where $\Aop=\Delta_F$ acts fiberwise (nonnegative, self-adjoint) and
\[
N: H^1(\Bint)\to L^2(\Bint)
\]
is locally Lipschitz on bounded sets (with Lipschitz constant $L$ on the relevant sublevel set).

The gradient flow is
\begin{equation}\label{eq:flow}
\partial_t\Psi(t) = -\Aop\Psi(t) - N(\Psi(t)),\qquad \Psi(0)=\Psi_0\in L^2(\Bint).
\end{equation}

\begin{remark}[Stabilization parameter]
Here $t$ parametrizes the analytic gradient/stabilization flow. It is not
identified with physical Lorentzian time. A physical evolution operator
$U(t_2,t_1)$ and its relation, if any, to this flow must be defined separately.
\end{remark}

\begin{assumption}[Well-posedness and smoothing]\label{ass:lwp}
For every $\Psi_0\in L^2$, \eqref{eq:flow} has a unique global mild solution
$\Psi\in C([0,\infty);L^2)\cap C((0,\infty);H^1)$.
For each $t>0$, the time-$t$ map $\Phi_t:\Psi_0\mapsto \Psi(t)$ maps $L^2$ into $H^1_F$
and is locally Lipschitz as a map $L^2\to H^1_F$; in particular, for some $M_1(t)$,
\[
\|\Phi_t(\Psi_1)-\Phi_t(\Psi_2)\|_{H^1_F}
\le M_1(t)e^{Lt}\|\Psi_1-\Psi_2\|_{L^2}.
\]
\end{assumption}

\begin{proposition}[Energy dissipation and forward invariance of sublevels]\label{prop:energydec}
Along sufficiently regular solutions of the $L^2$-gradient flow \eqref{eq:flow}, the energy is nonincreasing and
\[
C[\Psi(t)] + \int_{0}^{t}\|\nabla C(\Psi(s))\|_{L^2}^{2}\,\dd s = C[\Psi_0]\qquad (t\ge 0).
\]
In particular, every energy sublevel $\{\Psi:\ C[\Psi]\le M\}$ is forward invariant.
For semilinear analytic-semigroup flows this identity can be justified by standard approximation arguments; see, e.g., \cite{Amann1995,Henry1981}.
\end{proposition}
\newcommand{\DM}{D_M}
\[
\DM := \{\Psi\in H^1(\Bint):\ C[\Psi]\le M\}.
\]

\begin{assumption}[Forward invariance of energy sublevels]\label{ass:sublevel-inv}
For each $M>0$, the energy sublevel $\DM=\{\Psi\in H^1(\Bint): C[\Psi]\le M\}$ is forward invariant under the
flow $\Phi_t$ used in the Darbo argument (in particular, under the base-regularized flow if $\varepsilon>0$ is employed).
\end{assumption}


\begin{remark}[Semigroup/ODE splitting]
Decompose $L^2=\Ran(\Qop)\oplus \Ran(\PiCoh)$.
On $\Ran(\Qop)$ the sectorial operator $\Aop|_{\Ran(\Qop)}\ge \lamstar$ generates an analytic semigroup with the standard
$t^{-1/2}e^{-\lamstar t}$ bound for $\Aop^{1/2}e^{-t\Aop}$.
On $\Ran(\PiCoh)$ the flow reduces to an ODE $\partial_t u=-\PiCoh N(u)$ (when coherence invariance holds).
\end{remark}

\begin{assumption}[Coherence invariance]\label{ass:coh-inv}
We say coherence invariance holds if
\[
N(\Ran(\PiCoh))\subset \Ran(\PiCoh).
\]
Under this condition the coherent sector is invariant under the flow.
\end{assumption}

\subsection{Projected map and identification of equilibria}

Define
\[
T_t := \PiCoh\circ \Phi_t:\ L^2\to \Ran(\PiCoh).
\]

\begin{proposition}[Projected fixed points are equilibria]\label{prop:projfix}
Assume coherence invariance (Assumption~\ref{ass:coh-inv}).
If $\Psi^\ast\in \Ran(\PiCoh)$ satisfies $T_t(\Psi^\ast)=\Psi^\ast$ for some $t>0$, then
$\Phi_s(\Psi^\ast)=\Psi^\ast$ for all $s\ge 0$, and $\nabla C(\Psi^\ast)=0$.
\end{proposition}

\begin{proof}
Coherence invariance implies $\Phi_s(\Psi^\ast)\in\Ran(\PiCoh)$ for all
$s\ge0$, hence $T_t(\Psi^\ast)=\Phi_t(\Psi^\ast)=\Psi^\ast$. Applying the
energy identity of Proposition~\ref{prop:energydec} on $[0,t]$ gives
\[
0=C[\Phi_t(\Psi^\ast)]-C[\Psi^\ast]
=-\int_0^t\|\nabla C(\Phi_s(\Psi^\ast))\|_{L^2}^2\,\dd s.
\]
The integrand therefore vanishes, so the orbit is stationary and
$\nabla C(\Psi^\ast)=0$.
\end{proof}

\begin{remark}[Without coherence invariance]
If $N(\Ran\PiCoh)\not\subset \Ran\PiCoh$, then a fixed point of $T_t=\PiCoh\Phi_t$ need not be a true equilibrium of $\Phi_t$.
In that case $T_t$ can still be used to build invariant sets, but not to identify equilibria.
\end{remark}

\paragraph{What the projected-fixed-point result means.}
Projection can make a trajectory appear to return even while discarded
components continue to move.  Proposition~\ref{prop:projfix} rules out that
false return in the invariant coherent sector: because the full energy loses
exactly the squared gradient along the orbit, returning to the same coherent
state after positive flow time leaves no room for hidden dissipation.  This
is the logical bridge used by every existence route below.

\subsection{Detailed semigroup bound underpinning \FPpdf{$M_1(t)$}{M1(t)}}

\begin{lemma}[Semigroup bound underpinning $M_1(t)$]\label{lem:semigroup}
Let $\Qop:=\Id-\Pi_{\mathrm{coh}}$. Then for all $t>0$,
\[
\|A^{1/2}e^{-tA}\Qop\|_{L^2\to L^2}
\le C\,t^{-1/2}e^{-\lambda^\* t},
\qquad
\|e^{-tA}\Pi_{\mathrm{coh}}\|_{L^2\to L^2}=1.
\]
Consequently, for all $\Psi_1,\Psi_2\in L^2$,
\[
\|\Phi_t(\Psi_1)-\Phi_t(\Psi_2)\|_{H^1_F}
\le C\bigl(1+t^{-1/2}e^{-\lambda^\* t}\bigr)e^{Lt}\,
\|\Psi_1-\Psi_2\|_{L^2}.
\]
In particular, one may take
\[
M_1(t):=C\bigl(1+t^{-1/2}e^{-\lambda^\* t}\bigr).
\]
\end{lemma}


\begin{lemma}[Lipschitz continuity of $T_t$]\label{lem:Tlip}
For each $t>0$ there exists $K(t)$ such that for all $\Psi_{1,2}\in L^2$,
\[
\|T_t(\Psi_1)-T_t(\Psi_2)\|_{H^1_F}
\le C_\Pi K(t)e^{Lt}\|\Psi_1-\Psi_2\|_{L^2}.
\]
Moreover, \emph{compactness} of $T_t$ requires base regularization (see Section~\ref{sec:schauder}) or another mechanism that provides compactness in the base variable.
\end{lemma}

% ------------------------------------------------------------
\section{Existence of equilibria and projected fixed points}

There is no single universal compactness argument, so this section presents
three routes with different geometric costs.  The direct method minimizes the
energy on a weakly closed coherent sector.  Schauder uses genuine compactness
created by base regularization on a compact total space.  Darbo is designed
for a noncompact base: confinement makes the distant tail uniformly small,
while local smoothing supplies compactness on each bounded region.  A model
needs only one valid route, but it must state which route it uses.

\subsection{Variational existence on the coherent sector}

\begin{assumption}[Coercivity and weak lower semicontinuity]\label{ass:coercive}
$C:H^1(\Bint)\to\bbR$ is bounded below, coercive on $H^1(\Bint)$, and sequentially weakly lower semicontinuous.
\end{assumption}

\begin{theorem}[Variational existence]\label{thm:variational}
Assume coherence invariance (Assumption~\ref{ass:coh-inv}) and Assumption~\ref{ass:coercive}.
Then there exists $\Psi^\ast\in \Hcoh$ such that $\nabla C(\Psi^\ast)=0$.
In particular, $T_t(\Psi^\ast)=\Psi^\ast$ for all $t\ge 0$.
\end{theorem}

\subsection{Alternative route: Schauder (compact base) and base-regularization}\label{sec:schauder}

To obtain compactness of the time-$t$ map in the \emph{base} variable, introduce a base-regularized generator
\[
\Aeps := \Delta_F + \varepsilon \Delta_Y,\qquad \varepsilon\in(0,1],
\]
and consider the corresponding flow
\[
\partial_t\Psi = -\Aeps\Psi - N(\Psi).
\]

\begin{assumption}[Base-regularized well-posedness]\label{ass:base-reg}
For some $\varepsilon\in(0,1]$ and $\delta>0$, the flow with generator
$\Aeps$ is globally well-posed and, for every $t>0$, maps bounded subsets of
$L^2$ into bounded subsets of $H^{1+\delta}(\Bint)$. This is a standing
regularity hypothesis to be verified for each nonlinearity and model.
\end{assumption}

\begin{lemma}[Compactness for compact bases]\label{lem:compact}
If $Y$ is compact and Assumption~\ref{ass:base-reg} holds, then for each $t>0$,
$\Phi_t:L^2\to H^1$ is compact on bounded sets. Consequently
$T_t=\PiCoh\Phi_t$ is compact on bounded sets in $H^1$.
\end{lemma}

\begin{proof}
The total fibered space is compact when $Y$ and the fibers are compact. The
Rellich embedding $H^{1+\delta}\hookrightarrow\!\hookrightarrow H^1$ is then
compact. Compose this embedding with the smoothing in
Assumption~\ref{ass:base-reg} and the bounded map
$\PiCoh:H^1\to H^1$ from Lemma~\ref{lem:PiH1}.
\end{proof}

\begin{theorem}[Schauder fixed point]\label{thm:schauder}
Assume $Y$ is compact and Assumption~\ref{ass:base-reg} holds. Let $D\subset \Hcoh$ be nonempty, closed, bounded, and convex, and assume $T_t(D)\subset D$ for some $t>0$.
Then $T_t$ has a fixed point in $D$.
\end{theorem}

\subsection{Darbo condensing on noncompact bases}

\begin{definition}[Kuratowski measure of noncompactness]
For a bounded set $B$ in a metric space, $\alpha(B)$ is the infimum of $r>0$ such that $B$ can be covered by finitely many sets of diameter $<r$.
\end{definition}

\begin{assumption}[Confinement]\label{ass:confinement}
There exists a smooth $V:Y\to[0,\infty)$ with $V(y)\to\infty$ as $r(y):=\dist_Y(y,y_0)\to\infty$ for some basepoint $y_0\in Y$, and a constant $K\ge 0$ such that
\[
C[\Psi]\ge \int_Y V(y)\,\|\Psi(\cdot,y)\|_{L^2(F(y))}^2\,\dd y - K.
\]
\end{assumption}

\begin{lemma}[Uniform tail smallness]\label{lem:tails}
Under Assumption~\ref{ass:confinement}, for each $M>0$ and $\epsilon>0$ there exists $R>0$ such that
\[
\int_{\{r(y)>R\}} \|\Psi(\cdot,y)\|_{L^2(F(y))}^2\,\dd y \le \epsilon^2
\quad\text{whenever } \Psi\in H^1(\Bint)\ \text{and}\ C[\Psi]\le M.
\]
\end{lemma}



\begin{lemma}[Compact-local / small-tail decomposition]\label{lem:KR}
Assume Assumptions~\ref{ass:coercive}, \ref{ass:base-reg}, and
\ref{ass:confinement}, and fix $t>0$.
Let $\chi_R\in C_c^\infty(Y)$ satisfy $0\le \chi_R\le 1$, $\chi_R\equiv 1$ on $\{r\le R\}$ and
$\chi_R\equiv 0$ on $\{r\ge R+1\}$, and write $\pi:\Bint\to Y$ for the base projection.

Define the \emph{output cutoff} maps
\[
K_R(\Psi_0):=(\chi_R\circ\pi)\,\Phi_t(\Psi_0),
\qquad
S_R(\Psi_0):=(1-\chi_R\circ\pi)\,\Phi_t(\Psi_0).
\]
Then $\Phi_t=K_R+S_R$. Moreover:
\begin{enumerate}[label=\textup{(\roman*)}]
\item For each $M>0$, the map $K_R:\DM\to L^2(\Bint)$ is compact.
\item If $\DM$ is forward invariant under $\Phi_t$, then for each $M>0$,
\[
\sup_{\Psi_0\in \DM}\|S_R(\Psi_0)\|_{L^2}\to 0
\qquad\text{as }R\to\infty.
\]
\end{enumerate}
\end{lemma}


\begin{remark}[Why this fixes the ``tail linearity'' issue]
Tail smallness follows either from Proposition~3.2 for the original gradient flow,
or from Assumption~3.3 when working with the base--regularized flow used in the Darbo argument.


\end{remark}

\begin{proposition}[Condensing property on energy sublevels]\label{prop:condensing}
Assume Assumptions~\ref{ass:coercive}, \ref{ass:base-reg}, and
\ref{ass:confinement}, and fix $t>0$ and an energy level $M>0$.
On the invariant sublevel $D_M:=\{\Psi:\ C[\Psi]\le M\}$, the map $\Phi_t$ is \emph{Kuratowski-condensing} in the sense of Sadovski\u{\i}:
for every bounded set $B\subset D_M$ with $\alpha(B)>0$ one has $\alpha(\Phi_t(B))<\alpha(B)$.
Consequently, $T_t=\PiCoh\circ\Phi_t$ is condensing in $L^2$ on any
$T_t$-invariant subset of $D_M\cap \Ran(\PiCoh)$, because $\PiCoh$ is an
$L^2$-orthogonal projection and hence has norm one.
\end{proposition}


\begin{theorem}[Darbo--Sadovski\u{\i} fixed point]\label{thm:darbo}
Assume Assumptions~\ref{ass:coh-inv}, \ref{ass:coercive}, \ref{ass:base-reg}, and \ref{ass:confinement}.
Fix $t>0$ and an energy level $M>0$. Let
$D\subset D_M\cap\Ran(\PiCoh)$ be nonempty, $L^2$-closed, bounded, convex,
and invariant under $T_t$. Then $T_t:D\to D$ has a fixed point. By
Proposition~\ref{prop:projfix}, that fixed point is an equilibrium of the
gradient flow.
\end{theorem}

\paragraph{Interpretation of the three routes.}
The variational theorem establishes an equilibrium without first constructing
a compact time map.  Schauder establishes a fixed point of a compact map but
requires compact-base smoothing and an invariant convex set.  Darbo weakens
compactness to a strict reduction of noncompactness, at the cost of a
confining potential and invariant energy sublevel.  These are analytic
existence mechanisms; none says that the equilibrium is unique, dynamically
preferred, or physically realized.

% ------------------------------------------------------------
\section{Uniqueness of equilibria}

Existence answers whether at least one coherent stationary configuration is
available.  It does not answer whether two different initial configurations
can settle into different equilibria.  Strong monotonicity excludes that
possibility directly.  The inertial-manifold route instead first proves that
long-time motion is captured by finitely many modes and then asks for
monotonicity only on that reduced system.

\subsection{Route (A): strong convexity / strong monotonicity on the coherent sector}

\begin{assumption}[Strong monotonicity on $\Hcoh$]\label{ass:strongmono}
There exists $\mu>0$ such that for all $u,v\in \Hcoh$,
\[
\ip{\nabla C(u)-\nabla C(v)}{u-v}_{L^2}\ge \mu \|u-v\|_{L^2}^2.
\]
\end{assumption}

\begin{theorem}[Uniqueness on the coherent sector]\label{thm:uniqueH}
Under Assumptions~\ref{ass:coh-inv} and \ref{ass:strongmono}, $C|_{\Hcoh}$ has a unique minimizer $\Psi^\ast$,
which is the unique equilibrium in $\Hcoh$. Moreover,
\[
\|\Phi_t(u)-\Phi_t(v)\|_{L^2}\le e^{-\mu t}\|u-v\|_{L^2}\qquad \forall\,u,v\in\Hcoh.
\]
\end{theorem}

\subsection{Route (B): inertial manifold reduction}

Let $P_0$ be a finite-rank spectral-bundle projection and $Q_0:=\Id-P_0$.

\begin{assumption}[Spectral gap and Lipschitz nonlinearity]\label{ass:IMgap}
Assume there exists $\gamma>0$ such that
\[
\dist\big(\spec(\Aop|_{\Ran(Q_0)}),\,\spec(\Aop|_{\Ran(P_0)})\big)\ge \gamma,
\qquad
\|N(u)-N(v)\|_{L^2}\le L\|u-v\|_{L^2}.
\]
\end{assumption}

\begin{theorem}[Inertial manifold]\label{thm:IM}
If $L<\gamma/4$, there exists a finite-dimensional $C^1$ inertial manifold
$M=\{\xi+G(\xi):\ \xi\in\Ran(P_0)\}$ that is invariant and exponentially attracting.
\end{theorem}

\begin{remark}[References]
Inertial manifold constructions for semilinear parabolic equations under a spectral-gap condition are classical; see, for example, \cite{Temam1997,Henry1981}.
\end{remark}

\begin{assumption}[Reduced strong monotonicity]\label{ass:redmono}
There exists $\mu>0$ such that the reduced vector field on $\Ran(P_0)$ is strongly monotone with constant $\mu$.
\end{assumption}

\begin{theorem}[Uniqueness via inertial manifold]\label{thm:IMunique}
Under Assumptions~\ref{ass:IMgap} and \ref{ass:redmono}, the flow has a unique equilibrium $\Psi^\ast\in M$,
exponentially attracting on $M$.
If additionally coherence invariance holds and $\Ran(\PiCoh)\subset \Ran(P_0)$, then $\Psi^\ast\in \Hcoh$ and agrees with the coherent-sector variational minimizer.
\end{theorem}

\begin{example}[IM works while convexity fails]
Let $F=S^1_R$ and $N(u)=u^3-u$ (Allen--Cahn). $N$ is not convex so Assumption~\ref{ass:strongmono} fails.
On $\{\|u\|_{H^1}\le \rho\}$ one has $\|N'(u)\|\le 3\rho^2+1=:L(\rho)$, and $\gamma$ is comparable to the spectral gap $\lamstar=R^{-2}$.
Choosing $\rho$ so that $L(\rho)<\gamma/4$ yields the inertial manifold and uniqueness on it.
\end{example}

\paragraph{Practical distinction.}
The first route is stronger and cleaner when the full coherent energy is
uniformly convex.  The second can survive a nonconvex ambient energy, but its
conclusion is only as global as the attracting manifold and reduced
monotonicity hypotheses.  The Allen--Cahn example illustrates the intended
use: nonconvexity alone need not end the analysis, provided a verified
spectral separation controls the reduced dynamics.

% ------------------------------------------------------------
\section{Galerkin approximation for the steady state}

We next turn from existence to computation.  A finite Galerkin matrix is not
automatically an approximation of the abstract fiber flow: one must first
specify a uniformly elliptic steady operator and a finite-dimensional trial
space.  To ensure coercivity and uniqueness, the elliptic part includes a
strictly positive zeroth-order term, interpreted analytically as a mass or
confinement contribution.

\begin{assumption}[Coercive elliptic steady operator]\label{ass:coerciveAeps}
Let $\varepsilon>0$ and assume the steady operator has the form
\[
\widetilde A_\varepsilon := \Delta_F + \varepsilon\Delta_Y + W,
\]
where $W\ge w_0>0$ is a measurable multiplication operator and the associated
closed quadratic form has domain
\[
V_\varepsilon:=\{w\in H^1(\Bint):W^{1/2}w\in L^2\}.
\]
Examples include $W\equiv m_0^2$ on compact $Y$ and a confining potential
$W=V(y)\to\infty$ on a noncompact base.
\end{assumption}

We consider the steady equation
\begin{equation}\label{eq:steady}
\widetilde A_\varepsilon \Psi + N(\Psi)=0.
\end{equation}

\begin{definition}[Energy norm for $\widetilde A_\varepsilon$]
Under Assumption~\ref{ass:coerciveAeps}, define for $w\in V_\varepsilon$
\[
\|w\|_{\varepsilon}^2
:= \|\nabla_F w\|_{L^2}^2 + \varepsilon\|\nabla_Y w\|_{L^2}^2
+\int_{\Bint}W|w|^2.
\]
\end{definition}

\begin{lemma}[Strong monotonicity in the energy norm]\label{lem:strongmonoF}
Assume Assumption~\ref{ass:coerciveAeps}. Suppose $N$ is $L$-Lipschitz in $L^2$ on the set of interest and $L<w_0$.
Let $F:=\widetilde A_\varepsilon+N:V_\varepsilon\to
V_\varepsilon^\ast$. Then for all $u,v$ in that set,
\[
\langle F(u)-F(v),u-v\rangle_{V_\varepsilon^\ast,V_\varepsilon}
\ge \|u-v\|_\varepsilon^2-L\|u-v\|_{L^2}^2
\ge \mu_\varepsilon \|u-v\|_{\varepsilon}^2,
\]
where $\mu_\varepsilon := 1-L/w_0>0$.
\end{lemma}
\begin{theorem}[C\'ea-type quasi-optimal error in the energy norm]\label{thm:cea}
Assume Lemma~\ref{lem:strongmonoF}. Let $V_N\subset V_\varepsilon$ be a finite-dimensional Galerkin space and let
$\Psi^\ast$ solve $\widetilde A_\varepsilon\Psi + N(\Psi)=0$.
Let $\Psi_N\in V_N$ satisfy the Galerkin condition
\[
\langle \widetilde A_\varepsilon\Psi_N+N(\Psi_N),v
\rangle_{V_\varepsilon^\ast,V_\varepsilon}=0
\qquad\forall\,v\in V_N.
\]
Then $\Psi_N$ is unique (for each $N$) and
\[
\|\Psi^\ast-\Psi_N\|_{\varepsilon}
\le
\frac{1+L/w_0}{1-L/w_0}\,
\inf_{v\in V_N}\|\Psi^\ast-v\|_{\varepsilon}.
\]
If $W$ is also bounded above, the norm $\|\cdot\|_{\varepsilon}$ is
equivalent to $H^1$, hence this yields an $H^1$ error bound with an additional
explicit norm-equivalence constant. For a confining unbounded $W$, the result
is instead an estimate in the weighted form space $V_\varepsilon$.
\end{theorem}
\begin{proof}
The quadratic-form part has continuity constant one in
$\|\cdot\|_\varepsilon$. Moreover,
\[
|\langle N(u)-N(v),w\rangle|
\le L\|u-v\|_{L^2}\|w\|_{L^2}
\le \frac{L}{w_0}\|u-v\|_\varepsilon\|w\|_\varepsilon.
\]
Thus $F$ has continuity constant $1+L/w_0$ and strong-monotonicity
constant $1-L/w_0$. Galerkin orthogonality and the standard nonlinear
C\'ea argument give the displayed ratio. Strong monotonicity also gives
uniqueness of the finite-dimensional solution.
\end{proof}
\begin{remark}
In Theorem~6.4 the Lipschitz constant $L$ is taken with respect to the $L^2$ norm.
The conversion to the energy norm $\|\cdot\|_\varepsilon$ uses the inequality
\[
\|u\|_{L^2}\le w_0^{-1/2}\|u\|_\varepsilon,
\]
which explains the appearance of the factor $(1+L/w_0)$ in the C\'ea constant.
\end{remark}

\paragraph{What the error estimate does and does not certify.}
The theorem says that the computed solution is, up to the displayed stability
factor, as accurate as the best element available in the chosen Galerkin
space.  It therefore converts approximation quality into a rigorous solution
error.  It does not certify that the basis, coefficients, or operator were
selected by MTT geometry; those provenance questions must be settled before
this numerical theorem is invoked.


% ------------------------------------------------------------
\section{Examples and counterexamples}

\subsection{Compact fibers with fixed geometry: circles and tori have \FPpdf{$\lamone>0$}{lambda1>0}}
\begin{proposition}
If $F$ is compact, then the first positive Laplace eigenvalue $\lambda_1(\Delta_F)>0$.
\end{proposition}

\begin{corollary}[Circle/torus examples]
On $S^1_R$, $\spec(\Delta)=\{(k/R)^2:\ k\in\bbZ\}$, so $\lambda_1=R^{-2}>0$.
Similarly for flat tori $T^m$ one has $\lambda_1>0$.
\end{corollary}

\subsection{How the uniform gap \FPpdf{$\lamstar$}{lambda*} can fail}
If radii degenerate (e.g.\ $R(y)\to\infty$) or fibers become noncompact, then $\lamstar$ can drop to $0$, destroying uniform estimates.

\subsection{Effect of base potentials}
A smooth confining potential $V:Y\to[0,\infty)$ added to $C$ improves tail coercivity (Assumption~\ref{ass:confinement}) but does not change the fiber spectra.

\subsection{Worked example with explicit constants: \FPpdf{$Y=\Sone$}{Y=S1}, \FPpdf{$F(y)=\Sone_{R(y)}$}{F(y)=S1(R(y))}}
Assume $R\in C^1(S^1)$ with $0<R_{\min}\le R(y)\le R_{\max}$ and $|R'(y)|\le \rho$.
\begin{enumerate}[label=\textup{(\alph*)}]
\item \textbf{First eigenvalue:} $\lamstar=\inf_y \lambda_1(\Delta_{F(y)}) = 1/R_{\max}^2$.
\item \textbf{$H^1$ projector bound (one explicit model):}

with $\Aop(y)=\Delta_{F(y)}=-R(y)^{-2}\partial_\theta^2$ one computes
$\partial_y \Aop(y)= 2R'(y)R(y)^{-3}\partial_\theta^2$ and hence
$\|\partial_y \Aop(y)\|_{H^2\to L^2}\le 2\rho R_{\min}^{-3}$.
Using the Riesz formula and $\|(z-\Aop)^{-1}\|\le 2/\lamstar$ for $|z|=\lamstar/2$ yields a concrete
\[
\|\Pi\|_{H^1\to H^1}\le 1+\sup_y \|\partial_y\Pi(y)\|_{L^2\to L^2}
\le 1+\frac{4\rho R_{\max}^2}{R_{\min}^3}.
\]
\item \textbf{IM threshold:} taking $P_0=\Pi$ gives $\gamma=\lamstar=R_{\max}^{-2}$, and the IM criterion is $L<\gamma/4 = 1/(4R_{\max}^2)$.
\end{enumerate}

% ------------------------------------------------------------
\section{Auxiliary estimates for spectral-bundle projectors}

\begin{theorem}[Explicit $H^1_F$ bound for $\PiCoh$]\label{thm:PiBound}
Under Assumptions~\ref{ass:bounded-geom}--\ref{ass:gap} there exists
\[
C_\Pi \le 1 + \frac{C(K_0,\iota_0,m)}{\lamstar}
\]
such that $\|\PiCoh\|_{H^1_F\to H^1_F}\le C_\Pi$.
If the fibers are $y$-independent and flat, $C_\Pi=1$.
\end{theorem}

\begin{lemma}[Norm equivalence on the harmonic bundle]
On $\Ran(\PiCoh)$ one has $\nabla_F\Psi=0$ in the scalar connected-fiber case, hence
$\|\Psi\|_{H^1_F}=\|\Psi\|_{L^2}$. More generally, on $\Ran(\PiCoh)$ all Sobolev norms are uniformly equivalent by elliptic regularity on the harmonic bundle.
\end{lemma}

% ------------------------------------------------------------
\section{\FPpdf{\L{}ojasiewicz--Simon}{Lojasiewicz--Simon} convergence of the gradient flow}

The preceding sections can produce an equilibrium and, under stronger
hypotheses, uniqueness.  The next question is whether a full trajectory
actually converges rather than merely possessing stationary cluster points.
Analyticity and the \L{}ojasiewicz--Simon inequality prevent an orbit from
wandering indefinitely among nearby critical configurations.

\begin{assumption}[Analyticity and Palais--Smale]\label{ass:PS}
$C$ is real-analytic and satisfies a Palais--Smale condition on energy sublevels:
if $C[\Psi_n]$ is bounded and $\|\nabla C(\Psi_n)\|_{L^2}\to 0$, then $(\Psi_n)$ has a convergent subsequence in $H^1$.
We also assume the orbit under consideration remains in a bounded energy
sublevel on which the gradient flow is precompact in $H^1$ and the local
\L{}ojasiewicz--Simon inequality is valid near every cluster point.
\end{assumption}

\begin{theorem}[\L{}ojasiewicz--Simon convergence]
Assume Assumption~\ref{ass:PS}. For any solution $\Psi(t)$ of \eqref{eq:flow} with initial data in an energy sublevel,
there exists $\Psi_\infty$ with $\nabla C(\Psi_\infty)=0$ and
\[
\lim_{t\to\infty}\|\Psi(t)-\Psi_\infty\|_{H^1}=0.
\]
If coherence invariance holds and $\Psi(0)\in \Ran(\PiCoh)$, then $\Psi_\infty\in \Ran(\PiCoh)$.
\end{theorem}

\begin{remark}[Reference]
The underlying \L{}ojasiewicz--Simon inequality and convergence mechanism originate with \cite{Simon1983}.
\end{remark}

\begin{corollary}[Exponential rate under nondegeneracy]
If, in addition, the limiting equilibrium is nondegenerate in the relevant
phase space---equivalently, the local \L{}ojasiewicz--Simon exponent is
$1/2$, for example because the linearized Hessian has a positive spectral
gap after quotienting exact symmetries---then
$\|\Psi(t)-\Psi_\infty\|_{H^1}\le Ce^{-\eta t}$ for some $\eta>0$.
Uniqueness by itself does not imply an exponential rate.
\end{corollary}

% ------------------------------------------------------------
\section{Singular perturbation and the limit \FPpdf{$\epsdown$}{eps->0}}

Base regularization was introduced to obtain compactness, not as a claim that
the target model necessarily contains an $\varepsilon\Delta_Y$ term.  This
section checks when the artificial regularization can be removed without
losing the coherent equilibrium.  The extra condition on the incoherent
sector is essential: it is what prevents positive fiber modes from surviving
in the limit.

Let $\Psi_\varepsilon$ solve the base-regularized steady equation
\[
(\Delta_F+\varepsilon\Delta_Y)\Psi_\varepsilon + N(\Psi_\varepsilon)=0.
\]

\begin{assumption}[Uniform bounds]\label{ass:unifbds}
Assume coherence invariance and $\sup_{\varepsilon\in(0,1]}\|\Psi_\varepsilon\|_{H^1}<\infty$.
Assume either $Y$ is compact or a confining potential provides compactness on energy sublevels.
Assume also that $N$ is sequentially demicontinuous along this compactness
class: whenever $u_j\rightharpoonup u$ in $H^1$ and $u_j\to u$ in $L^2$,
then $N(u_j)\rightharpoonup N(u)$ in $H^{-1}$. This holds for the standard
subcritical Nemytskii nonlinearities but must be checked in each model.
\end{assumption}

\begin{theorem}[Convergence as $\varepsilon\downarrow 0$]
Under Assumption~\ref{ass:unifbds} and, in addition, either
\begin{enumerate}[label=\textup{(\roman*)}]
\item the $\Qop$-sector Lipschitz constant satisfies $L<\lamstar$, or
\item $N$ is weakly monotone on the $\Qop$-sector:
$\ip{\Qop N(u)-\Qop N(v)}{\Qop(u-v)}_{L^2}\ge 0$ for all $u,v$,
\end{enumerate}
there exists a subsequence $\varepsilon_j\downarrow 0$ with $\Psi_{\varepsilon_j}\to \Psi^\ast$ in $L^2$ and weakly in $H^1$,
where
\[
\Delta_F\Psi^\ast + N(\Psi^\ast)=0,\qquad \text{and necessarily }\Psi^\ast\in \Ran(\PiCoh).
\]
\end{theorem}

\begin{proof}
The uniform $H^1$ bound and the compactness hypothesis give a subsequence with
the stated weak-$H^1$/strong-$L^2$ convergence. The term
$\varepsilon_j\Delta_Y\Psi_{\varepsilon_j}$ tends to zero in $H^{-1}$,
while the fiber form passes to the limit weakly and demicontinuity identifies
the nonlinear limit. This yields the limiting steady equation.

Let $p=\PiCoh\Psi^\ast$ and $q=\Qop\Psi^\ast$. Coherence invariance gives
$\Qop N(p)=0$. Testing the $\Qop$ equation against $q$ gives
\[
\lamstar\|q\|_{L^2}^2
\le -\langle \Qop(N(p+q)-N(p)),q\rangle.
\]
Under alternative (i), the right side has absolute value at most
$L\|q\|_{L^2}^2$, so $L<\lamstar$ forces $q=0$. Under alternative (ii),
the nonlinear term is nonnegative on the left-hand side and again $q=0$.
Thus $\Psi^\ast\in\Ran(\PiCoh)$.
\end{proof}

% ------------------------------------------------------------
\section{Related work and novelty}

Our existence and regularity statements are classical in spirit, but three aspects appear new or at least not treated in this precise combination:
(i) spectral-bundle projections with uniform constants on fibered products with bounded geometry (Theorem~\ref{thm:PiBound}), including explicit constants in a worked model;
(ii) a condensing proof on noncompact bases via a compact-local/small-tail decomposition (Lemma~\ref{lem:KR}) tied directly to confinement;
(iii) a clean link between projected fixed points and true equilibria (Proposition~\ref{prop:projfix}) via coherence invariance.

\section{Conclusion and boundary of the result}

The paper constructs a conditional chain from bundle geometry to a computable
coherent equilibrium.  A uniform fiber gap defines the harmonic projector;
bounded geometry controls that projector over the base; dissipative flow and
coherence invariance identify projected returns with true equilibria;
variational, compact, or condensing arguments establish existence; additional
monotonicity or spectral separation supplies uniqueness; and the coercive
Galerkin theorem provides a quantitative approximation once the steady
operator is fixed.

The principal achievement is therefore not the unconditional existence of an
MTT vacuum.  It is a reusable checklist of analytic gates with explicit
failure modes.  Loss of the uniform gap invalidates global projector
estimates.  Lack of coherence invariance permits false projected returns.
Absence of base compactness or confinement defeats the corresponding
fixed-point route.  Existence without monotonicity leaves multiple equilibria
possible, and a discretization without operator provenance remains only a
numerical model.

The next Fixed Points papers may build on this analytic foundation to study
projected equilibria, damping balances, curvature coupling, and physical
interpretation.  The separate foundational task is to intertwine the abstract
coherent sector used here with the selected q79 vertical geometry and its
shared-circle structure.  Until that realization theorem is supplied, every
physical application of the present results remains conditional on the
displayed hypotheses.

% ------------------------------------------------------------
\bibliographystyle{plain}
\bibliography{refs}

\end{document}
